\documentclass[10pt,twocolumn]{IEEEtran}

\usepackage{amsmath, amssymb, amsthm, bm, mathtools, array}
\usepackage{graphicx}
\usepackage{algorithmic}
\usepackage{balance}

\newtheorem{theorem}{Theorem}
\newtheorem{lemma}{Lemma}
\newtheorem{definition}{Definition}
\newtheorem{corollary}{Corollary}

\DeclareMathOperator*{\argmin}{arg\,min}
\DeclareMathOperator*{\argmax}{arg\,max}
\DeclareMathOperator{\diag}{diag}
\DeclareMathOperator{\rank}{rank}
\DeclareMathOperator{\prox}{prox}
\DeclareMathOperator{\Tr}{Tr}
\DeclareMathOperator{\KL}{KL}
\DeclareMathOperator{\det}{det}
\DeclareMathOperator{\MSE}{MSE}
\DeclareMathOperator{\Var}{Var}
\DeclareMathOperator{\E}{\mathbb{E}}
\DeclareMathOperator{\P}{\mathbb{P}}

\begin{document}
	
	\title{Prepare for the Worst, Hope for the Best: Active Robust Learning On Distributions}
	\author{Seyed Hossein Ghafarian and Hadi Sadoghi Yazdi}
	
	\maketitle
	
	\begin{abstract}
		In recent years, many learning systems have been developed for higher level forms of data, such as learning on distributions in which each example itself is a distribution. This article proposes active robust learning on distributions. In learning on distributions, there is no access to distributions themselves but rather access is through a sample drawn from a distribution. Therefore, similar to robust learning, any estimates of examples are inexact. In order to address these difficulties, we provide an upper bound on the risk of the classifier in the next stage of active learning, where the size of the labeled dataset increases. Based on this upper bound, we propose probabilistic minimax active learning (PMAL) as a general multiclass active learning method that is easy to use in many Bayesian settings, which provably selects an example with knowledge of its label minimizing the expected risk. We present an efficient approximation of the objective with a known error bound to deal with the intractability of the proposed method for active robust learning. Here, we face a nonconvex problem, which we solve by means of a related convex problem with a bound on the norm of the difference between their solutions. To utilize the information about the estimates of distributions, we propose active robust learning on the distributions method based on learning the kernel embedding of distributions by a recent Bayesian method. The experiments demonstrate the effectiveness of the resulting method on a set of synthetic and real-world distributional datasets.
	\end{abstract}
	
	\begin{IEEEkeywords}
		Active learning, kernel embedding of distributions, learning on distributions.
	\end{IEEEkeywords}
	
	\section{Introduction}
	
	Learning systems require knowledge to successfully learn a concept. The goal of learning systems is to predict the label of unseen test data. Without knowledge, learning systems are unable to generalize the concept of learning to the unseen. In learning systems, both a priori knowledge or human user knowledge can be utilized.
	
	Knowledge transfer from humans to machines is a fundamental aspect of learning systems. Active learning is a paradigm for the transfer of knowledge from humans to machines, in which the machine asks for the supervision of some selected pieces of data. In order to minimize the supervisory efforts needed, active learning aims at intelligently using limited supervision resources and human knowledge, by smartly selecting useful pieces of data and asking human users to supervise them. However, humans usually view knowledge in more abstract or higher level forms, while algorithms and machines usually require lower level forms of knowledge, such as labels of simple vectors.
	
	In practice, knowledge transfer using higher level forms of data can be very advantageous in active learning, since users may be able to more easily provide information. In these situations, the active learning algorithm may acquire richer information from the user. For example, Du and Ling \cite{du2010asking} proposed a "generalized" query, equivalent to many specific queries, for a rich interaction between humans and machines.
	
	Recently, learning systems have become increasingly more complex, operating on higher level forms of data, such as structured or distributional data. In learning from structured data, data have some structures, such as graphs \cite{getoor2003learning}; whereas, in learning on distributions \cite{muandet2012learning}, each data point itself is a distribution. It is highly desirable to build effective methods for minimizing the supervision of these forms of data.
	
	In many applications, it is more natural to assume that knowledge is not in the instances themselves, but in the distribution from which the instances are drawn. For example, in tasks, such as anomaly detection, each individual point may be normal, but a group of points may be anomalous \cite{poczos2012support}. In these applications, it is better to directly model the problem as learning on distributions. In recent years, many methods for learning on distributions have been developed \cite{xiong2014learning, oliva2014fast, muandet2017kernel, szabo2016learning}.
	
	The current article proposes an approach for active learning on distributions. In learning on distributions, examples have inherent uncertainty. Each distributional example is estimated by employing a sample with a limited size and, therefore, has uncertainty which is a challenge in active learning.
	
	To address this challenge in a principled manner, we first consider the active learning problem in the most general setting. We provide an upper bound on the risk of the classifier in the next stage of active learning with an increased size of the labeled dataset even though we have not yet acquired any new example label. An example is queried such that its selection minimizes this upper bound. We call the resulting method probabilistic minimax active learning (PMAL). PMAL can be used in many Bayesian settings.
	
	Based on PMAL, an active robust learning algorithm is proposed. Here, examples can be inexact. The proposed learning algorithm is robust against the inexactness in examples. This is even more challenging in active learning when there is only partial information about data. Unfortunately, the computation of PMAL for active robust learning is intractable. To deal with intractability, we present an efficient approximation for probabilistic minimax active robust learning (PMARL).
	
	In this approximation technique, we encounter a nonconvex optimization problem. We found an upper bound on the norm of the difference between a regularized solution of this problem and a convex problem. The approximation is tight in the sense that some labels exist for unlabeled data, where the approximation error is negligible.
	
	The problem of learning on distributions may be viewed as a special case of robust learning. While we do propose a general active robust learning algorithm, which considers uncertainty in the estimation of examples, our focus is on how to deal with uncertainty when estimating distributional examples. Furthermore, in this case, robust learning differs from some other research \cite{ghafarian2016active, xu2009robustness}, whose goal is to be robust against outliers. To the best of our knowledge, this is the first active learning method specially designed for learning on distributions.
	
	Any learning on distributions algorithm relies on estimating distributions. Recently, some approaches for better estimation of distributions in kernel methods have been proposed \cite{muandet2016kernel, muandet2017kernel}. In Section I-A2 we see that these estimates are uncertain.
	
	Many of the algorithms for learning on distribution ignore inherent estimation uncertainty. Employing a recent method known as Bayesian learning of kernel embedding \cite{flaxman2016bayesian}, the present article presents a method for active robust learning on distributions. It is called PMARL on distributions which utilizes Bayesian learning of kernel embedding (PMARLDB) to directly deal with the uncertainty in the estimation of kernel mean embeddings. Experiments demonstrate that the proposed method is very effective in utilizing this inherent uncertainty and in selecting examples of distributions using samples.
	
	To sum up, three main contributions are made by the current work.
	\begin{enumerate}
		\item PMAL that is general and can be applied to many Bayesian settings. This method has been proven to select examples, which minimize the expected risk of the classifier.
		\item Active robust learning and the proposed approximation method that deals with its associated nonconvex problem with a known error bound.
		\item An algorithm for active learning on distributions, which is robust in the sense that the estimates of each example feature uncertainty.
	\end{enumerate}
	
	In the next section, we first review related works.
	
	\subsection{Related Work}
	
	\subsubsection{Active Learning}
	Most of the active learning methods try to minimize the number of queries by selecting the most representative and/or informative examples.
	
	Informative methods view the problem of selecting examples from the perspective of parameter space, that is, classifier space, while representative methods see the problem from the perspective of input space. Our method, on the other hand, concentrates on a joint conditional probability of labels given examples and ask label of points with the maximum probability of misclassification. Therefore, the proposed method is both representative and informative.
	
	Querying the label of the most representative examples reduces the labeling effort, since they make it easier for an algorithm to more accurately infer many other labels from the structure of the data. In \cite{dasgupta2008hierarchical}, a hierarchical clustering approach for measuring representativeness is employed. Cai and He \cite{cai2012manifold} proposed a manifold optimal experiment design method. The proposed method considers representativeness based on the usefulness of unlabeled examples in estimating a closer conditional probability of labels given examples to its true value. The main drawback drawback of representative approaches is that their performance heavily depends on the cluster or manifold assumption of the data. Furthermore, the main drawback of representative methods is that they ignore the actual label structure of data.
	
	By considering the label structure of the data, informative methods usually avoid this drawback. The most informative examples are those which more quickly minimize the size of the set of classifiers capable of separating classes in data, that is, version space. In the proposed method, we prove that the example selected for the query is the best example in that it minimizes the expected risk of the classifier.
	
	Uncertainty sampling \cite{lewis1995sequential} is an informative approach that selects instances with the highest uncertainty in the current model. Zhang et al. \cite{zhang2015active} tackled the difficult problem of imbalanced and label noisy data in active learning for crowdsourcing by uncertainty sampling. Another major informative approach is query by committee, which selects an example with most disagreement on its label among a set of classifiers trained on the labeled data \cite{freund1997selective}. Inspired by query by committee, Wang et al. \cite{wang2017committee} proposed an evolutionary PSO optimization algorithm.
	
	Active learning for support vector machines (SVMs) is introduced in \cite{tong2001support}. This method queries the label of the example $x^* = \argmin_{x \in D_U} |f(x)|$, where $D_U$ is unlabeled data and $f(x)$ is the SVM. This objective minimizes version space by selecting the nearest example to the current decision boundary \cite{tong2001support}. A min--max objective is proposed to minimize function $f$ \cite{hoi2008semi}. In \cite{hoi2008semi}, the representativeness of examples is considered through manifold learning. Huang et al. \cite{huang2014active} viewed this as a heuristic approach and replaces it with minimizing the labels of unlabeled examples. However, Huang et al. \cite{huang2014active} failed to provide any theoretical justification for its approach. The present work differs from \cite{hoi2008semi} and \cite{huang2014active} in that we introduce PMAL which is more general and provably consistent in that its expected risk will converge to Bayes risk with an increasing number of labeled instances.
	
	The study by Du et al. \cite{du2017exploring} introduces the usage of a two-sample discrepancy problem and an uncertainty measure for representativeness and informativeness. Du et al. \cite{du2017exploring} estimated the probability of each instance given the model using the LIBSVM output \cite{chang2011libsvm} for estimating the similarity between labeled and unlabeled data using two-sample discrepancy measure in a quadratic objective function. The current work differs in that we directly use a probabilistic formulation.
	
	Many active learning methods are heuristic. Fang et al. \cite{fang2017active} proposed a principled selection method that minimizes the upper bound of the risk in a multitask setting with trace norm-regularized least squares, which take advantage of jointly learning multiple tasks and asking queries. Our method is based on minimizing a different upper bound of the risk of the classifier in a single task setting via a completely different approach. Ramirez-Loaiza et al. \cite{ramirez2017active} provided a comprehensive analysis of basic methods. For a survey on active learning methods, see \cite{settles2012active}.
	
	\subsubsection{Learning on Distributions}
	Many natural phenomena produce distributions as their output. Therefore, it may be much more advantageous to directly model them as distributions rather than to transform them into a vector, as does most research in machine learning. In learning on distributions, each example is a distribution. The objective is a learning algorithm, for example, a classifier that can learn the label of an input distribution.
	
	This concept became popular with support measure machines (SMMs) \cite{muandet2012learning}. SMM solves an SVM problem on distributional data with the expected kernel. Letting $P_i$ and $P_j$ be two distributions, then the expected kernel for $P_i$ and $P_j$ is
	\begin{equation}
		\langle K_{P_i, P_j} \rangle = \langle \mu^*_{P_i}, \mu^*_{P_j} \rangle = \mathbb{E}_{x_i \sim P_i, x_j \sim P_j} \left[ k(x_i, x_j) \right].
	\end{equation}
	
	Kernel methods developed for learning on distributions usually employ kernel mean embedding of distributions, $\mu^*_{P_i} := \int_{\mathcal{Z}} k(z, \cdot) dP_i(z)$, in which a distribution $P_i$ is represented as a mean function in a Hilbert space associated with kernel function $k$ \cite{smola2007hilbert}. $\mu^*_{P_i}$ is unique if $k$ is a characteristic kernel \cite{muandet2017kernel}.
	
	The challenge is that there is no access to a distribution itself, but rather to a sample drawn from it. Thus, it is a two-stage problem in which estimates have to rely on similarities computed between two sets of points and not directly between two distributions \cite{szabo2016learning}. Using samples $\{x_{iu}\}_{u=1}^{n_i}$ and $\{x_{jv}\}_{v=1}^{n_j}$, the empirical estimate of the expected kernel is usually employed instead
	\begin{equation}
		\langle K_{P_i, P_j} \rangle = \frac{1}{n_i n_j} \sum_{u=1}^{n_i} \sum_{v=1}^{n_j} k(x_{iu}, x_{jv}).
	\end{equation}
	
	Since the sample size for a distribution is limited, empirical kernel mean embedding of distribution $P_i$, $\hat{\mu}_{P_i}$ (Table I) is employed, which is just a mean of the kernel function for every instance in the sample \cite{muandet2017kernel}. Recently, it has been shown that empirical kernel mean embedding is not the best estimate of distributions and some methods for obtaining better estimates have been proposed \cite{muandet2016kernel, flaxman2016bayesian}. Let $\mathcal{H}$ be the Hilbert space associated with kernel function $k$. For any $\mu_P \in \mathcal{H}$, $\mu_P(x)$ means evaluating $\mu_P$ at location $x$.
	
	Recently, Szabo et al. \cite{szabo2016learning} have researched the learning-theoretic analysis of distribution regression. Interestingly, Lopez-Paz et al. \cite{lopez2015towards} studied casual inference as a distribution classification problem. Moreover, Lopez-Paz et al. \cite{lopez2015towards} explored the theoretical properties of both distribution classification and casual inference.
	
	Sample sizes drawn from distributions often vary greatly. We use the following theorem in Section II-B to devise a simple active robust learning. Lopez-Paz et al.'s \cite{lopez2015towards} theorem provides the convergence of $\hat{\mu}_{P_i}$, the empirical mean embedding of distribution $P_i$, through a sample size of $n_i$, to its population counterpart $\mu^*_{P_i}$, in the reproducing Kernel Hilbert space (RKHS) $\mathcal{H}_k$.
	
	\begin{theorem} \label{thm:empirical-convergence}
		Assume that $\|f\|_\infty \le 1$, for all $f \in \mathcal{H}_k$ with $\|f\|_{\mathcal{H}_k} \le 1$. Then, with the probability at least $1 - \delta$
		\begin{equation}
			\|\mu^*_{P_i} - \hat{\mu}_{P_i}\|_{\mathcal{H}_k} \le \frac{2\sqrt{\mathbb{E}_{z \sim P_i}[k(z, z)]}}{\sqrt{n_i}} + \sqrt{\frac{2 \log \frac{1}{\delta}}{n_i}}.
		\end{equation}
	\end{theorem}
	
	Based on the above theorem, the uncertainty of estimates and the similarities between distributions are dependent on the sample size. However, much research employing kernel methods does not consider the dependence of uncertainty on the sample sizes of distributions. Law et al. \cite{law2017bayesian} considered this uncertainty without computing the uncertainty in the kernel Hilbert space and proposed an RBF network formulation for regression using a set of landmark points.
	
	In the proposed active learning on distributions method, we take into account both the uncertainty in estimating distributions and the dependence of this uncertainty on sample sizes. For a more detailed review of learning on distributions, refer to \cite{szabo2016learning}. Muandet et al. \cite{muandet2017kernel} provided an excellent comprehensive review on kernel mean embedding, including learning on distributions.
	
	Distributional data pose a set of challenges for active learning algorithms. In the next section, we propose a method to effectively face these challenges.
	
	\section{Proposed Method}
	
	Examples are shown in lowercase letters and bold lowercase letters indicate the ordered list of examples, such as $\mathbf{x} = (x_1, \ldots, x_n)$ or their labels $\mathbf{y} = (y_1, \ldots, y_n)$. We denote the labeled subset of data with $L$ and unlabeled subset with $\bar{U}$. Assume a subset $Q \subset \bar{U}$ with size $b$ and let $U = \bar{U} - Q$. For brevity, we denote $L \cup Q$ with $L_Q$. For a set $S \subseteq \{1, \ldots, n\}$ with length $s$, we denote $\mathbf{x}_S = (x_{S(1)}, x_{S(2)}, \ldots, x_{S(s)})$. Thus, $\mathbf{x}_L$ or $\mathbf{y}_{L_Q}$ means labeled examples and labels of labeled plus query examples, respectively. For brevity, notations are provided in Table I. Also, proofs of lemmas and theorems are in appendices.
	
	As mentioned earlier, there are no exact estimates of distributions. This setting is similar to robust learning, but it also differs in that there is specific information about estimates of distributions using samples, which can also be utilized.
	
	To effectively utilize this information, we first propose the PMAL method in Section II-A. Using PMAL, we deal with inexact examples in active robust learning, in a principled manner. Section II-B proposes an algorithm for active robust learning in input space. It also provides the basics for active robust learning on distributions. Furthermore, Section II-B introduces a simpler form of active robust learning on distributions. Section II-C, proposes the main objective of the current paper, which is active robust learning on distributions. In the following, some necessary definitions are discussed.
	
	\subsection{Definitions and Background}
	Consider the random pair $(X, Y)$, which "is defined by the pair" $(\mu, \eta^*)$, "where $\mu$ is the probability measure for $X$" and $\eta^*$ "is the regression of $Y$ on $X$" \cite{devroye1996probabilistic}. Then, the (unknown) conditional probability of labels given examples is $\eta^*(x, y) = \P(Y = y|X = x)$.
	
	\begin{definition}[Bayes Decision Function and Bayes Risk]
		The Bayes risk and Bayes classifier are, respectively, defined as $L^* = \P(g^*(X) \neq Y)$ and $g^*(x) = \argmax_{y \in \mathcal{C}}\{\eta^*(x, y)\}$ \cite{devroye1996probabilistic, chen2006consistency}.
	\end{definition}
	
	\begin{definition}
		Let $y$ be the corresponding label for $x$. Consider classifier function $g_D(x)$, also known as the plug-in decision function trained using set $D$. Its corresponding expected risk is $L_D = \P\{g_D(X) \neq Y|D\}$ \cite{devroye1996probabilistic}. It is easy to see that
		\begin{align}
			L_D &= 1 - \P\{g_D(X) = Y|D\} \\
			&= 1 - \sum_{y \in \mathcal{C}} \P\{g_D(X) = y, Y = y|D\} \\
			&= 1 - \sum_{y \in \mathcal{C}} \int_{\mathbb{R}^d} \mathbb{I}_{[g_D(x)=y]} \P\{Y = y|D, X = x\} \mu(dx).
		\end{align}
		We define $\P\{Y = y|D, X = x\}$ as $\eta$, $\eta_D$ interchangeably
		\begin{equation}
			\eta(x, y; \mathbf{x}_D, \mathbf{y}_D) \equiv \eta_D(x, y) := \P\{Y = y|D, X = x\}.
		\end{equation}
		Also, let $A(x, D) := \max_{y \in \mathcal{C}}\{\eta^*(x, y) - \eta(x, y; \mathbf{x}_D, \mathbf{y}_D)\}$. To minimize the risk, the plug-in classifier $g_D$ is equal to $g_D(x) = \argmax_{y \in \mathcal{C}}\{\eta(x, y; \mathbf{x}_D, \mathbf{y}_D)\}$ \cite{chen2006consistency}.
	\end{definition}
	
	In other words, to minimize the risk the non-negative function $\eta$ must approximately be equal to $\eta^*$ at $(x, y)$ by using data, $(\mathbf{x}, \mathbf{y})$, as parameters. In fact, it is essential to have some assumptions, which relates $\mu$, $\eta^*$ and $\eta$, that is, "the process generating data and the process of assigning labels" \cite{sinha2009semi} to be able to make effective use of unlabeled examples \cite{singh2009unlabeled}. Tsybakov uses the following assumption which holds for "various types of estimators and various sets of probability distributions" $\mathcal{P}$, \cite{devroye1996probabilistic, audibert2007fast}:
	
	\begin{assumption} \label{ass:tsybakov}
		Consider $\eta_D$ as an estimator of the regression function $\eta^*$ using data $D$ and $\mathcal{P}$, a class of probability distributions on $(X, Y)$ such that for some constants $C_1 > 0$ and $C_2 > 0$, and a positive sequence $a_n$, for almost all $x$ with respect to $P_X$, we have
		\begin{equation}
			\sup_{P \in \mathcal{P}} \P\{A(x, D) \ge \delta\} \le C_1 \exp(-C_2 a_n \delta^2).
		\end{equation}
	\end{assumption}
	
	\begin{lemma} \label{lem:unlabeled-benefit}
		Assume unlabeled data help in reducing risk of the classifier on $\mathcal{P}$, under Assumption \ref{ass:tsybakov}, then $\exists c_U > 1$, such that with probability $1 - C_1 \exp(-C_2 a_n \delta^2)$
		\begin{equation}
			\int_{\mathbb{R}^d} A(X, D) \mu(dx) \le c_U^{-1} \int_{\mathbb{R}^d} A(X, L) \mu(dx) + 2\delta.
		\end{equation}
	\end{lemma}
	
	\subsection{Probabilistic Minimax Active Learning}
	
	In each stage of active learning, we would like to select an example that if its label is known, will help build an accurate model. Our best hope is that the current model correctly predicts almost all the labels of unlabeled data, which we call the best labels for the current model or, simply, the best labels. However, the worst case is when the probability is high for a reverse label of an example. Here, a reverse label refers to a label that is different from a label predicted by the current model.
	
	Since there is no information about the label of an unlabeled example until we query the user, it is necessary to be prepared for the worst case regarding the label of the query; at the same time, we must be hopeful that the labels of other unlabeled examples have the best label. Note that we prefer a simpler model. Therefore, we expect the best label for unlabeled examples, except for the queried label, because we assume that the current model is simpler than models in the next stages of active learning where the number of labeled examples increases. Hence, prepare for the worst and hope for the best.
	
	In a stage of active learning, only subset $D_L$ is used for training and the labels of unlabeled instances $y_i, i \in U$ are unknown. We ask what will happen to classifier accuracy if we add an example to the labeled data before we even know its label. Theorem \ref{thm:risk-upper-bound} provides an upper bound on the risk of a classifier trained using labeled data plus the unlabeled example(s) $Q$. First, we prove Theorem \ref{thm:plug-in-risk-bound} for multiclass plug-in classifiers, which, to the best of our knowledge, did not appear before in the literature.
	
	\begin{theorem} \label{thm:plug-in-risk-bound}
		Let $T_D(x, y) = \eta^*(x, y) - \eta_D(x, y)$. Then
		\begin{equation}
			L_D - L^* \le \sum_{i=1}^{|\mathcal{C}|} \mathbb{E}|T_D(X, c_i)| \le |\mathcal{C}| \max_{y \in \mathcal{C}} \mathbb{E}|T_D(X, y)|.
		\end{equation}
	\end{theorem}
	
	\begin{theorem} \label{thm:risk-upper-bound}
		Assume that assumptions of Lemma \ref{lem:unlabeled-benefit} hold. Let $\Delta_L \equiv \zeta \int_{\mathbb{R}^d} \max_{y \in \mathcal{C}}\{\eta^*(x, y)-\eta(x, y; \mathbf{x}, \mathbf{y})\}\mu(dx)$ and $\beta = 1-c_U^{-1}$. Then, with probability $1 - C_1 \exp(-C_2 a_n \delta^2)$, the upper bound for the risk of a classifier trained using the label of the query $Q$, $L_{L_Q}$, is
		\begin{equation}
			L_{L_Q} - L^* \le \beta \Delta_L + 2\zeta \delta + \zeta \alpha \left( 1 - c_0 \min_{\mathbf{y}_Q} \max_{\mathbf{y}_U} P(\mathbf{y}|\mathbf{x}) \right) E_\eta,
		\end{equation}
		where $E_\eta = \mathbb{E}(\max_{y \in \mathcal{C}} \eta(X, y; \mathbf{x}_L, \mathbf{y}_L))$, and $\alpha > 0$, $c_0 > 1$, and $\zeta \ge 2$ are some constants.
	\end{theorem}
	
	The question is how to know which examples label will improve classification accuracy the most. The answer is the examples that minimize the upper bound on the expected risk for the classifier trained using labeled data plus the query.
	
	\begin{corollary} \label{cor:pmal-query}
		Based on Theorem \ref{thm:risk-upper-bound}, the best query, which minimizes the upper bound on the expected risk, is
		\begin{equation}
			Q = \argmax_{\{Q, |Q|=b\}} \min_{\mathbf{y}_Q} \max_{\mathbf{y}_U} P(\mathbf{y}|\mathbf{x}).
		\end{equation}
	\end{corollary}
	
	This theorem essentially states that, in order to achieve higher accuracy, we need to query labels of examples which, regardless of the learner, most improve our knowledge of the conditional probability of labels given examples, $P(Y|X)$. If we know this distribution better, it will be easier to build a learner for it. We know that $\forall \mathbf{y}_U, \forall \mathbf{y}_Q$
	\begin{equation}
		\min_{\mathbf{y}_Q} P(\mathbf{y}|\mathbf{x}) \le \min_{\mathbf{y}_Q} \max_{\mathbf{y}_U} P(\mathbf{y}|\mathbf{x}) \le \max_{\mathbf{y}_U} P(\mathbf{y}|\mathbf{x}).
	\end{equation}
	
	At the same time, we know that when using all available knowledge, there are still some examples whose class probability is still low. This occurs in regions, where we expect misclassification to occur \cite{rasmussen2004gaussian}, leading to objective (5). In this case, to achieve higher accuracy levels, we must concentrate on examples/regions, whose probability of their reverse labels is the highest using all available knowledge. Comparing (5) with (6), reveals that selected example bisects the unlabeled data.
	
	If the size of batch ($b$) is small, objective (5) can be solved easily. However, if it is large, approximations such as relaxation \cite{goemans1995improved} or heuristic techniques can be used. We leave it for future works.
	
	Our main goal is to propose a method for active learning on distributional data. As seen later, there is access to each distribution $P_i$ through a sample $\{x_{i1}, x_{i2}, \ldots, x_{in}\}$. Therefore, we do not have an accurate estimate of $P_i$. In order to develop a method for active learning on distributions, this inaccuracy must be considered when selecting examples. For this reason, the next section proposes a method for active learning on input space data with measurement error.
	
	\subsection{Bayesian Active Robust Learning for Input Space}
	
	In a robust learning setting, each example is inaccurate. There may be many reasons for the inaccuracy in examples, such as measurement error. In any case in a robust learning setting, we assume there is a probability of error in examples. On the other hand, the performance of the learning algorithm must not degrade due to errors. In other words, the algorithm must be robust against these inaccuracies. Let $x_i$ denote an instance and $x_i^*$ denote its true value. Here, we assume a linear model and that each data point $x_i$ has an uncertainty represented by the following conditional distribution:
	\begin{equation}
		y_i = x_i^{*T} w + \epsilon, \quad x_i^*|x_i \sim \mathcal{N}(0, \sigma_i^2), \quad \epsilon \sim \mathcal{N}(0, \sigma^2).
	\end{equation}
	
	Then, in order to compute $P(\mathbf{y}|\mathbf{x})$, we first compute $P(\mathbf{y}|\mathbf{x}, w) = \prod_i P(y_i|x_i, w)$. There is
	\begin{align}
		P(y_i|x_i, w) &= \int P(y_i|w, x_i^*) P(x_i^*|x_i) dx_i^* \\
		&= \int \mathcal{N}(y_i; x_i^{*T} w, \sigma^2) \mathcal{N}(x_i^*; x_i, \sigma_i^2) dx_i^* \\
		&= \mathcal{N}(y_i; x_i^T w, \sigma^2 + \sigma_i^2 \|w\|^2).
	\end{align}
	The first equality is the total law of probability. The second equality comes from (7). The last equation is the well-known integration of two Gaussians. Let $\sigma_i(w) = \sigma^2 + \sigma_i^2 \|w\|^2$. We have
	\begin{equation}
		P(\mathbf{y}|\mathbf{x}, w) \propto \prod_i \sigma_i(w)^{-\frac12} \exp\left( -\frac12 \sum_i \frac{(y_i - w^T x_i)^2}{\sigma_i(w)} \right).
	\end{equation}
	
	In other words, if a point has a large variance, it must have less influence on the classifier and, therefore, the probability of selecting this point must be small. This also shows that if a classifier is simple, that is, if $\|w\|$ is small, the variance of the conditional probability $P(\mathbf{y}|\mathbf{x}, w)$, will be less affected by an example's uncertainty.
	
	Let us assume a prior $w \sim \mathcal{N}(0, \sigma_w^{-2} I)$. The integral
	\begin{equation}
		P(\mathbf{y}|\mathbf{x}) = \int P(\mathbf{y}|\mathbf{x}, w) P(w) dw
	\end{equation}
	is intractable. More importantly, $\mathbf{y}_U$ is unknown. Therefore, we have to find an approximation of $P(\mathbf{y}|\mathbf{x})$ as a function of $\mathbf{y}_U$. We assume a Gaussian probability density function (pdf) for the approximation $\tilde{P}(\mathbf{y}|\mathbf{x}, w)$
	\begin{equation}
		\tilde{P}(\mathbf{y}|\mathbf{x}, w) = \prod_{i=1}^n \mathcal{N}(y_i; w^T x_i, \sigma^2 + \sigma_i^2 v^2)
	\end{equation}
	where $v$ will be revealed a little later. Using $\tilde{P}$ instead of $P$ solves our problem of intractability of the integral. But how to know how much they are close for every value of $\mathbf{y}_U$ at the situation, which we cannot compute the integral?
	
	The following theorem approximates KL divergence between $P$ and $\tilde{P}$. It is based on recent results in Laplace approximation of intractable integrals \cite{wong2001asymptotic, lapinski2019multivariate, inglot2014simple}.
	
	\begin{theorem} \label{thm:kl-approx}
		Let $w^* = \argmax_w P(\mathbf{y}, w|\mathbf{x})$ and $\lambda = 1 / \max_i \sigma_i^2$. Also, let
		\begin{align}
			h(w; \mathbf{y}, \mathbf{x}) &= \frac12 \sum_{i=1}^n \left( \frac{(y_i - w^T x_i)^2}{\sigma_i(w)} + \log \sigma_i(w) \right) + \frac{1}{2\sigma_w^2} \|w\|^2, \\
			g(w, v) &= \frac12 \sum_{i=1}^n \left( \frac{(\sigma_i^2 \|w\|^2 - v^2)(y_i - w^T x_i)^2}{\sigma_i(v) \sigma_i(w)} + \log \frac{\sigma_i(v)}{\sigma_i(w)} \right).
		\end{align}
		Then, we have
		\begin{equation}
			\KL(P(\mathbf{y}, w|\mathbf{x}) || \tilde{P}(\mathbf{y}, w|\mathbf{x})) = c_\lambda \left\{ \frac{n D(P, \tilde{P}; \mathbf{y}, \mathbf{x}, \lambda, v)}{\sqrt{\det D^2 h(w^*; \mathbf{y}, \mathbf{x})}} + \epsilon(\lambda) \right\},
		\end{equation}
		where
		\begin{equation}
			D(P, \tilde{P}; \mathbf{y}, \mathbf{x}, \lambda, v) \propto e^{-h(w^*;\mathbf{y},\mathbf{x})} \frac{g(w^*, v)}{\sqrt{\det D^2 h(w^*; \mathbf{y}, \mathbf{x})}},
		\end{equation}
		$c_\lambda = (2\pi/\lambda)^{d/2}$ and $\epsilon(\lambda) = e^{-h(w^*;\mathbf{y},\mathbf{x})}(O(1)/\sqrt{\lambda})$ when $\lambda \to \infty$.
	\end{theorem}
	
	If we set $v = \|w^*\|$ based on Theorem \ref{thm:kl-approx}, $g(w^*, v)$ becomes $0$ and, therefore, the KL divergence between $\tilde{P}$ and $P$ will be nearly $0$.
	
	However, since $h$ is nonconvex, computing $w^*$ is difficult. The following lemma provides an upper bound for $h$. This shows that instead of directly computing $w^*$, we can minimize a simple convex function in which the function $w$ is more regularized. Interestingly, Lemma \ref{lem:h-upper-bound} relates robust learning to (more) regularization of the learning function. This has been shown in an SVM robust learning setting in \cite{xu2009robustness}.
	
	\begin{lemma} \label{lem:h-upper-bound}
		Let $\sigma_w^{-2} = \sigma^2 ( \frac{1}{\sigma_w^2} + \sum_{i=1}^n \frac{\sigma_i^2}{\sigma^2} )$ and
		\begin{equation}
			h'(w; \mathbf{y}, \mathbf{x}) = \frac12 \sum_{i=1}^n \frac{(y_i - w^T x_i)^2}{\sigma^2} + \frac12 \frac{\sigma_w^{-2}}{\sigma^2} \|w\|^2.
		\end{equation}
		Then, $h(w; \mathbf{y}, \mathbf{x}) \le h'(w; \mathbf{y}, \mathbf{x}) + n \log \sigma$.
	\end{lemma}
	
	Therefore, $h'$ can be used to minimize $h$. However, it is important to know the difference of their solutions to use $h'$ instead of $h$.
	
	\begin{theorem} \label{thm:w-diff-bound}
		Let $w^*$ be a regularized answer of the problem $\argmin_w h(w; \mathbf{y}, \mathbf{x})$ and $\tilde{w} = \argmin_w h'(w, \mathbf{y}, \mathbf{x})$. Then
		\begin{equation}
			\frac{\|w^* - \tilde{w}\|}{\|w^*\|} \le \sigma_w^2 \sigma_w^2 \left( D_x + N_x \right) \|w^*\| + \sigma_w^2 \left( \sum_{i=1}^n \frac{\sigma_i^2}{\sigma^2} \right),
		\end{equation}
		$\lim_{\sigma_1, \ldots, \sigma_n \to 0} \|w^* - \tilde{w}\| = 0$ and
		\begin{align}
			D_x &= \sum_{i=1}^n \sum_{j=1}^n \frac{\|D_{ij}\| |\sigma_j^2 - \sigma_i^2|}{\sigma_i(w^*) \sigma_j(w^*)} \le \sum_{i=1}^n \sum_{j=1}^n \frac{\|D_{ij}\|}{\sigma^4} |\sigma_j^2 - \sigma_i^2|,
		\end{align}
		where $D_{ij} = \sigma^{-2}(x_j x_j^T + \sigma_j^2 I)(x_i y_i)$ and
		\begin{equation}
			N_x = \frac{1}{\sigma_w^2 \sigma^2} \sum_{i=1}^n \frac{\|x_i y_i\| \sigma_i^2}{\sigma_i(w^*)} \le \frac{1}{\sigma_w^2 \sigma^4} \sum_{i=1}^n \|x_i y_i\| \sigma_i^2.
		\end{equation}
		Note that $\lim_{\|w^*\| \to \infty} \frac{D_x}{\|w^*\|^2} = \lim_{\|w^*\| \to \infty} N_x = 0$.
	\end{theorem}
	
	Since $\tilde{w}$ is sufficiently close to $w^*$, we can use it as an approximation to $w^*$ for computing $\tilde{P}$ and based on the above theorem and Theorem \ref{thm:kl-approx}, the distance between $P$ and $\tilde{P}$ is sufficiently small for any value of $\mathbf{y}_U$. Since $e^{-h(w^*;\mathbf{y},\mathbf{x})}$ and $e^{-h'(w^*;\mathbf{y},\mathbf{x})}$ are changing much faster with respect to $\mathbf{y}_U$ than $g$, we use $\min_{\mathbf{y}_U} \min_w h'(w; \mathbf{y}, \mathbf{x})$ for the approximation.
	
	Let $K$ be the kernel matrix and $L = (K + \sigma_w^{-2} I)^{-1} K (K + \sigma_w^{-2} I)^{-1}$.
	
	\begin{theorem} \label{thm:p-approx}
		Let $\tau = \mathbf{y}_L^T (L_{ll} - L_{lu} L_{uu}^{-1} L_{ul}) \mathbf{y}_L$, $\kappa_i^{-1} = \sigma^2 + \sigma_i^2 \tau$, $\Upsilon = \diag(\kappa)$, and $S = \Upsilon - \Upsilon \{ \sigma_w^2 (K^{-1} + \sigma_w^2 \Upsilon)^{-1} \} \Upsilon$, we have
		\begin{equation}
			\tilde{P}(\mathbf{y}|\mathbf{x}) \propto \exp\left( -\frac12 \mathbf{y}^T S \mathbf{y} \right).
		\end{equation}
	\end{theorem}
	
	The final objective for PMARL is
	\begin{equation}
		Q = \argmin_{\{Q, |Q|=b\}} \max_{\mathbf{y}_Q} \min_{\mathbf{y}_U} \mathbf{y}^T S \mathbf{y}.
	\end{equation}
	
	Given the quadratic objective, it is easy to compute $Q$ with $b = 1$. There are many real-world applications of active robust learning. One is active learning on distributions, as presented in the next section. Objective (14) can be used for active learning on distributions. In order to do so, we can set $\sigma_i = \beta n_i^{-1/2}$ and $\Upsilon = \diag(1 \oslash (\sigma^2 + \sigma_I \tau))$, where $\beta$ is a parameter. This method is called PMARL on distributions using similarity (PMARLDS). Section III-D compares this method with that of the next section.
	
	\subsection{Probabilistic Minimax Active Robust Learning On Distributions Using Bayesian Learning of Kernel Embedding}
	
	To make a clear distinction, the point on which learning is performed is called an \emph{example}, while a point in a sample drawn from a distribution is called an \emph{instance}. There are a number of distributions $P_i, i \in \{1, \ldots, n\}$ as our examples. There is no direct access to a distribution. Rather, from each distribution $P_i$, we have drawn a sample $\mathbf{x}_i = \{x_{i1}, \ldots, x_{i n_i}\}$, where each instance in the sample is a vector, that is, $x_{iu} \in \mathbb{R}^d, \forall u \in [1, n_i]$. We need to compute $P(\mathbf{y}|\hat{\boldsymbol{\mu}}_P)$.
	
	Embedding $\hat{\mu}_{P_i}$ is not the best estimate for distribution $P_i$ as shown by Muandet et al. \cite{muandet2017kernel}. Recently, it has been shown that it is possible to estimate a more accurate kernel embedding with an uncertainty measure in the Hilbert space in a Bayesian framework \cite{flaxman2016bayesian}. We use this framework to consider the uncertainty in estimating kernel mean embedding for distributions. Letting $k_\theta$ be a positive-definite kernel with parameter $\theta$, Bayesian kernel embedding (BKE) \cite{flaxman2016bayesian} proposes employing a prior on $\mu^*$ as follows so as to ensure that $\mu^* \in \mathcal{H}_k$:
	\begin{align}
		\mu^*|\theta &\sim \mathcal{GP}(0, r_\theta(\cdot, \cdot)), \\
		r_\theta(x, y) &= \int k_\theta(x, u) k_\theta(u, y) du,
	\end{align}
	where $\mathcal{GP}$ is a Gaussian process. The choice of $r_\theta$ ensures that $\mu^*_{P_i} \in \mathcal{H}$ \cite{flaxman2016bayesian}. $\theta$ is the parameters of the kernel function and from now on, it will not be shown in the equations.
	
	The following lemma gives the probability of a label of example $P_i$, given the evaluation of kernel mean embedding of example $P_i$ at each point $x_{it} \in \mathbf{x}_i, t \in [1, n_i]$.
	
	\begin{lemma} \label{lem:prob-label-embedding}
		Let $r$ be the kernel function as in (15), $r_{xx} = r(x, x)$, $R_i^x = [r(x, x_{i1}), r(x, x_{i2}), \ldots, r(x, x_{i n_i})]^T$ and $R_i = [r(x_{it}, x_{iv})]_{t,v=1}^{n_i}$. Also, let $B_i = (R_i + (\rho^2 / n_i) I_n)^{-1} R_i^x$ and $C_i = r_{xx} - R_i^{xT} (R_i + (\rho^2 / n_i) I_n)^{-1} R_i^x$. Let $\hat{\mu}_{B_i}(x_i) = B_i \hat{\mu}_{P_i}(x_i)$. Assuming model $y_i = w^T \mu^*_{P_i} + \epsilon, \epsilon \sim \mathcal{N}(0, \sigma^2)$, there is
		\begin{equation}
			P\left(y_i|\hat{\mu}_{P_i}(x_i), w\right) \propto \sigma_i(w)^{-\frac12} \exp\left( -\frac{(y_i - w^T \hat{\mu}_{B_i}(x_i))^2}{2 \sigma_i(w)} \right),
		\end{equation}
		where $\sigma_i(w) = \sigma^2 + w^T C_i w$.
	\end{lemma}
	
	If distribution $P_i$ has a large covariance $C_i$ in the sense that $w^T C_i w$ becomes large, then its effect on the classifier will be little. Therefore, the probability of selecting $P_i$ for query in active learning must be low. By looking at $C_i$ in Lemma \ref{lem:prob-label-embedding}, we can gain insight onto when this will happen. This occurs when normalized $R_i^x$, that is, $(R_i + (\rho/n_i) I_n)^{-1/2} R_i^x$, is near perpendicular to $w$ in the Hilbert space.
	
	Similar to Section II-B, it is intractable to compute the integral $P(\mathbf{y}|\hat{\boldsymbol{\mu}}_P(\mathbf{x})) = \int P(\mathbf{y}|\hat{\boldsymbol{\mu}}_P(\mathbf{x}), w) P(w) dw$. We use the same technique as in Section II-B. Therefore, similarly, we use a pdf $\tilde{P}$ instead of $P$ in the above integral and follow the exact same method.
	
	\begin{theorem} \label{thm:p-approx-dist}
		Let $\mu_{R_i} = (R_i + (\rho^2 / n_i) I_n)^{-1} \hat{\mu}_{P_i}(x_i)$, and let $K_{\mu B}$ be a matrix with elements $K_{\mu B}(i, j) = [\mu_{R_i}^T R_{ij} \mu_{R_j}]$. Also, let
		\begin{equation}
			P_t = [R_{ti} \mu_{R_i}]_{i=1}^{n} \in \mathbb{R}^{n_t \times n}, \quad E_t = K_{\mu B} - P_t^T (R_t + (\rho^2 / n_t) I_n)^{-1} P_t, \quad L_t = (K_{\mu B} + \sigma_w^{-2} I)^{-1} E_t (K_{\mu B} + \sigma_w^{-2} I)^{-1}.
		\end{equation}
		Moreover, $\tau_i = \mathbf{y}_L^T \{ L_i^{ll} - L_i^{ul^T} (L_i^{uu})^{-1} L_i^{ul} \} \mathbf{y}_L$, $\kappa_i^{-1} = \sigma^2 + \tau_i$, and $\Upsilon = \diag(\kappa)$. Then
		\begin{equation}
			\tilde{P}(\mathbf{y}|\hat{\boldsymbol{\mu}}_P(\mathbf{x})) \propto \exp\left( -\frac12 \mathbf{y}^T S \mathbf{y} \right),
		\end{equation}
		where
		\begin{equation}
			S = \Upsilon - \sigma_w^2 \Upsilon \{ K_{\mu B} - K_{\mu B} \{ \sigma_w^{-2} \Upsilon^{-1} + K_{\mu B} \}^{-1} K_{\mu B} \}.
		\end{equation}
	\end{theorem}
	
	This is very intuitive, since it essentially describes a similarity between distributions $(i, j)$, for selecting distribution $t$'s label, which is based on the similarity of embedding their instances in their distribution with the common ground, that is, distribution $t$'s metric.
	
	\section{Experimental Results}
	
	This section provides a detailed experiment of active learning on distributions. First, to show the effectiveness of the currently proposed method in a controlled setting, we perform a set of experiments on a synthetic dataset. After that, a set of experiments is performed on real-world datasets. In each experiment, we compare our method with the following methods.
	\begin{enumerate}
		\item \textbf{DISTRAND}: This method randomly selects distributional examples.
		\item \textbf{DISTQUIRE}: We computed the means of each distribution in input space. Then, for queries, we select examples using the QUIRE method \cite{huang2014active}.
		\item \textbf{DISTHARM}: The means of each distribution are computed and then examples are selected by the harmonic Gaussian method \cite{zhu2003semi}.
		\item \textbf{DISTMARGIN}: Using the computed means of each distribution, margin-based active learning selects informative examples \cite{tong2001support}.
		\item \textbf{HILBQUIRE}: Here, we propose the HILBQUIRE method, which selects examples using the Quire method \cite{huang2014active}, in which the kernel matrix between examples in QUIRE is computed by using mean embedding of distributions.
		\item \textbf{PMARLDB}: This is our proposed active robust learning on distribution method.
	\end{enumerate}
	
	In all the experiments, each dataset is randomly divided into two disjoint datasets for training and testing data. At each stage in active learning, we train a classifier on labeled data and report accuracy on test data. We start active learning by selecting one example randomly for each class. Parameter $\rho$ in Lemma \ref{lem:prob-label-embedding} is set to $\{0.1, 1\}$. Also, the median rule is used for the kernel bandwidth. When datasets are large, for example, Reuters, we subsampled $20\%$ of examples. We implement classifier SMM \cite{muandet2012learning} by using LIBSVM \cite{chang2011libsvm} for training and testing on distributional data. We use an empirical expected kernel with RBF kernel between distributions using (1).
	
	\subsection{Synthetic Dataset}
	
	The synthetic distributional dataset in \cite{muandet2012learning} is employed and we call it GAUSSDIST. Two classes of Gaussian distributions are considered to be positive and negative classes. The mean parameters of these distributions are normally distributed around $(1, 1, \ldots, 1)$ and $(2, 2, \ldots, 2)$ with a covariance of $\Sigma = 0.5 I$. The covariance of these distributions is generated using the Wishart distributions with a covariance of $\Sigma_+ = 0.6 I$ and $\Sigma_- = 1.2 I$ with 10 degrees of freedom.
	
	We select $30\%$ of the learning set for the test set. The learning set consists of 150 distributions for the positive and negative classes. We employ the input space kernel parameter, level-2 kernel parameters to $\gamma = \gamma_{is} = 0.1$, and regularization parameter $\sigma_w^{-2} = 2^{-7}$. Each experiment is repeated 50 times to compensate for the randomness.
	
	\begin{figure}[h]
		\centering
		\includegraphics[width=0.45\textwidth]{fig1a.png}
		\includegraphics[width=0.45\textwidth]{fig1b.png}
		\caption{Comparison of classification accuracy for dataset GAUSSDIST with (a) different sample sizes (b) equal sample sizes for each of distributions. (a) $\{5, 100\}$ instances for each distribution. (b) 5 instances for each distribution.}
		\label{fig:gaussdist}
	\end{figure}
	
	In order to show the effectiveness of the currently proposed active learning method, which considers uncertainty in estimating kernel mean embedding of distributions, we design an experiment with comparably different sample sizes. To achieve this, sample sizes of 5 or 100 are randomly used. The results of the experiments are presented in Fig. 1. When the numbers of instances drawn from each distribution differ, the performance of our PMARLDB method is superior to that of the other methods, as some distributions with many instances may bias the active learning method.
	
	By directly modeling the uncertainty in estimating distributions, the current method's accuracy increases when uncertainty is high. Our approach uses the uncertainty information in both high and low sample size distributions to select better distributions. Even when the number of instances from each distribution is equal and low, our method performs better than do other methods. Although HILBQUIRE performs well in a different sample size scenario, it does not do well in low sample sizes. This usually happens when there is much uncertainty in estimating the kernel mean embedding of distributions used in HILBQUIRE.
	
	\subsection{Real-World Datasets}
	
	\subsubsection{Handwritten Character Recognition}
	Here, results are presented from experiments performed on the handwritten digit dataset, USPSDIST \cite{muandet2012learning}. This dataset is a distributional dataset created from images of digits 0--9. Instances are produced by basic transformation (scaling, translation, and rotation) on images. Each image of a handwritten digit is a $16 \times 16$ matrix. A Gaussian prior over parameters of transformation determines the equivalent classes of distributions.
	
	There are two pairs of parameters for scaling and translation, one for the $x$-axis and the other for the $y$-axis. For rotation, we have a single parameter as the angle of rotation. Similar to \cite{muandet2012learning}, $\mathcal{N}([1, 1], 0.1 I_2)$, $\mathcal{N}([0, 0], 5 I_2)$, and $\mathcal{N}(0, \pi)$ are utilized for scaling, translation and rotation, respectively. We randomly select a subset of pairs of digits, 6 versus 9, 3 versus 8, 3 versus 4, and 1 versus 8. In each of the classification tasks mentioned above, 100 distribution points are used. For each distribution point $P_i$, a random number of instances $n_i$ are generated by the same method as in \cite{muandet2012learning}. The number $n_i$ is set to $\max(1, \text{round}(\mathcal{N}(2, 10)))$.
	
	The active learning experiments start with two labeled examples and finish with 50 labeled distributional examples. Experiments for each dataset were repeated 50 times to compensate for randomness. The results are shown in Fig. 2.
	
	\begin{figure}[h]
		\centering
		\includegraphics[width=0.45\textwidth]{fig2a.png}
		\includegraphics[width=0.45\textwidth]{fig2b.png}
		\includegraphics[width=0.45\textwidth]{fig2c.png}
		\includegraphics[width=0.45\textwidth]{fig2d.png}
		\caption{Comparison of classification accuracy for the USPDSDIST dataset. (a) Digits 6 and 9. (b) Digits 3 and 8. (c) Digits 3 and 4. (d) Digits 1 and 8.}
		\label{fig:uspsdist}
	\end{figure}
	
	As seen in Fig. 2, even with a small variance in the sample sizes of distribution points, the current study's proposed PMARLDB method outperforms other methods in almost all cases. A comparison of our approach, PMARLDB, with random sampling reveals that PMARLDB selects distributional points whose labels contribute much more to the performance of the classifier.
	
	\subsubsection{Bag of Words}
	This section presents experiments on some text classification datasets. For the representation of documents, the bag of words representation is utilized. We assume that each document can be represented as a distribution from which words have been drawn. Each word can be represented as a vector. Since documents consist of a variable number of words, there are sets of different cardinalities representing documents.
	
	Here, word2vec \cite{mikolov2013efficient} is employed, which is quite effective for the representation of words. For the experiments, we use the 20newsgroup dataset$^1$ which consists of documents in 20 classes. Since it is a multiclass dataset, we selected a number of more difficult binary tasks and a few simple ones.
	
	Because some of the documents have a very large number of words in the 20newsgroup dataset, documents with more than 100 words were removed for these experiments. We randomly selected 300 documents for each class. Fig. 3 presents the experimental results. As seen in Fig. 3, the proposed method, PMARLDB is superior to almost all others and consistently outperforms random sampling. For difficult tasks, such as "alt.atheism" and "talk.religion.misc," our approach is quite successful in selecting distributions.
	
	\begin{figure}[h]
		\centering
		\includegraphics[width=0.45\textwidth]{fig3a.png}
		\includegraphics[width=0.45\textwidth]{fig3b.png}
		\includegraphics[width=0.45\textwidth]{fig3c.png}
		\includegraphics[width=0.45\textwidth]{fig3d.png}
		\includegraphics[width=0.45\textwidth]{fig3e.png}
		\includegraphics[width=0.45\textwidth]{fig3f.png}
		\includegraphics[width=0.45\textwidth]{fig3g.png}
		\includegraphics[width=0.45\textwidth]{fig3h.png}
		\includegraphics[width=0.45\textwidth]{fig3i.png}
		\includegraphics[width=0.45\textwidth]{fig3j.png}
		\includegraphics[width=0.45\textwidth]{fig3k.png}
		\includegraphics[width=0.45\textwidth]{fig3l.png}
		\includegraphics[width=0.45\textwidth]{fig3m.png}
		\includegraphics[width=0.45\textwidth]{fig3n.png}
		\includegraphics[width=0.45\textwidth]{fig3o.png}
		\includegraphics[width=0.45\textwidth]{fig3p.png}
		\caption{Comparison of classification accuracy for the 20newsgroup dataset. Each of the figures above shows a binary classification between two classes shown by class numbers. Class numbers 1--20 are is shown in Table II. (a) 1 and 18. (b) 1 and 20. (c) 2 and 4. (d) 2 and 6. (e) 3 and 4. (f) 4 and 7. (g) 5 and 6. (h) 5 and 7. (i) 5 and 13. (j) 8 and 11. (k) 9 and 10. (l) 10 and 11. (m) 12 and 15. (n) 16 and 17. (o) 16 and 18. (p) 19 and 20.}
		\label{fig:20news}
	\end{figure}
	
	\subsection{Multiclass Experiments}
	
	In this section, we present the result of experiments on multiclass distributional datasets. To perform multiclass experiments, we use the simplest scheme to represent classes. We assign an odd integer (negative or positive) as labels. We used this scheme for multiclass experiments of our method. More sophisticated handling of class labels may produce better results. We use the following multiclass datasets.
	
	\begin{enumerate}
		\item \textbf{COIL Unprocessed and COIL Processed Datasets:} These datasets consist of pictures of 20 objects from $0^\circ$ to $360^\circ$, with $5^\circ$ increments \cite{nene1996columbia}. These are seamlessly fit into the notion of distributional example. We consider a randomly chosen set of pictures of each object from different angles as a distributional example.
		\item \textbf{Article Sentences Dataset:} This dataset is made from scientific articles from PLoS Computational Biology, the machine-learning repository on arXiv, and the psychology journal Judgment and Decision Making. argumentative zones (AZ), which is a scheme for classifying the function of sentences, is used which includes five labels \cite{teufel2002summarizing}.
		\item \textbf{Reuters-21578:} This dataset is a well-known dataset. It is a set of stories from Reuters in 1987. The documents are in eight categories.
	\end{enumerate}
	
	Except for binary methods (MARGIN and QUIRE), we also considered multiclass uncertainty sampling plus other methods for experiments as in the previous experiments. The results of experiments can be seen in Fig. 4. As it can be seen in Fig. 4, the proposed method has also a good performance in multiclass active learning.
	
	\begin{figure}[h]
		\centering
		\includegraphics[width=0.45\textwidth]{fig4a.png}
		\includegraphics[width=0.45\textwidth]{fig4b.png}
		\includegraphics[width=0.45\textwidth]{fig4c.png}
		\includegraphics[width=0.45\textwidth]{fig4d.png}
		\caption{Comparison of classification accuracy for multiclass datasets. (a) Processed COIL. (b) Unprocessed COIL. (c) Article Sentence. (d) Reuters.}
		\label{fig:multiclass}
	\end{figure}
	
	\subsection{Analysis on the Effect of Bayesian Learning of Kernel Embedding}
	
	A question that naturally comes to mind is if the proposed active learning on distributions method, PMARLDB, is more effective in selecting examples and achieving a better accuracy through the Bayesian learning of kernel embeddings.
	
	Alternatively, the active robust learning of Section II-B can also tackle the problem of active learning on distributions. In order to characterize the uncertainty in examples, there is Theorem \ref{thm:empirical-convergence}, which describes the uncertainty for empirical kernel mean embedding of distributions. This uncertainty can be considered as a Gaussian distribution with a mean of empirical kernel mean embedding and a variance proportional to $(1/\sqrt{n})$. It is essentially a degraded version of our PMARLDB method in that the Bayesian learning of kernel mean embedding is not considered. In other words, we must ask if the active robust learning method of Section II-B can be effective on distributions employing Theorem \ref{thm:empirical-convergence} for uncertainty in examples.
	
	Theorem \ref{thm:empirical-convergence} essentially states that the estimate of distributions using the kernel mean embedding of distributions has a variance of order $O(n^{-1/2})$. Therefore, one can use the formulation of Section II-B to select examples. As discussed earlier, we call this method PMARLDS. This section compares PMARLDS with the currently proposed main method so as to evaluate the latter's effectiveness. To achieve this, we perform experiments on the 20newsgroup dataset to compare PMARLDB with PMARLDS. The results are shown in Fig. 5 and indicate that PMARLDB always outperforms PMARLDS. PMARLDB is superior due to simultaneously learning of BKE of distributions.
	
	\begin{figure}[h]
		\centering
		\includegraphics[width=0.45\textwidth]{fig5a.png}
		\includegraphics[width=0.45\textwidth]{fig5b.png}
		\includegraphics[width=0.45\textwidth]{fig5c.png}
		\includegraphics[width=0.45\textwidth]{fig5d.png}
		\caption{Comparison of classification accuracy between Bayesian Active Learning in Hilbert (PMARLDS) and PMARLDB for the 20newsgroup dataset. Each of the figures above shows a binary classification between two classes shown by class numbers. Class numbers 1--20 are shown in Table II. (a) 1 and 2. (b) 2 and 13. (c) 4 and 8. (d) 4 and 15.}
		\label{fig:pmarlds_vs_pmarldb}
	\end{figure}
	
	\section{Conclusion}
	
	In this article, we proposed a general and easy-to-use active learning method, called PMAL, based on the minimization of the risk upper bound. We used PMAL to propose a method for active robust learning in which examples are inexact. This method has been used as a basis to propose an active learning on distributions method. In the future, we plan to develop techniques for improving the scalability of the proposed method. It is also interesting to extend PMAL to other learning tasks.
	
	\balance
	
	\bibliographystyle{IEEEtran}
	\begin{thebibliography}{1}
		
		\bibitem{du2010asking}
		J. Du and C. X. Ling, ``Asking generalized queries to domain experts to improve learning,'' \emph{IEEE Trans. Knowl. Data Eng.}, vol. 22, no. 6, pp. 812--825, Jun. 2010.
		
		\bibitem{getoor2003learning}
		L. Getoor, N. Friedman, D. Koller, and B. Taskar, ``Learning probabilistic models of link structure,'' \emph{J. Mach. Learn. Res.}, vol. 3, pp. 679--707, Mar. 2003.
		
		\bibitem{muandet2012learning}
		K. Muandet, K. Fukumizu, F. Dinuzzo, and B. Schölkopf, ``Learning from distributions via support measure machines,'' in \emph{Proc. Neural Inf. Process. Syst.}, 2012, pp. 10--18.
		
		\bibitem{poczos2012support}
		B. Póczos, L. Xiong, D. J. Sutherland, and J. Schneider, ``Support distribution machines,'' 2012. [Online]. Available: papers2://publication/uuid/BFE768AB-7981-4607-AE79-2EAF3C5DC299
		
		\bibitem{xiong2014learning}
		L. Xiong and J. G. Schneider, ``Learning from point sets with observational bias,'' in \emph{Proc. Uncertainty Artif. Intell.}, 2014, pp. 898--906.
		
		\bibitem{oliva2014fast}
		J. B. Oliva, W. Neiswanger, B. Poczos, J. Schneider, and E. Xing, ``Fast distribution to real regression,'' in \emph{Proc. AISTATS}, vol. 33, 2014, pp. 706--714.
		
		\bibitem{muandet2017kernel}
		K. Muandet, K. Fukumizu, B. Sriperumbudur, and B. Schölkopf, ``Kernel mean embedding of distributions: A review and beyond,'' \emph{Found. Trends Mach. Learn.}, vol. 10, nos. 1--2, pp. 1--141, 2017.
		
		\bibitem{szabo2016learning}
		Z. Szabo, B. Sriperumbudur, B. Poczos, and A. Gretton, ``Learning theory for distribution regression,'' \emph{J. Mach. Learn. Res.}, vol. 17, no. 1, pp. 5272--5311, Nov. 2016.
		
		\bibitem{ghafarian2016active}
		H. Ghafarian and H. S. Yazdi, ``Active robust learning,'' 2016. [Online]. Available: arXiv:1608.07159.
		
		\bibitem{xu2009robustness}
		H. Xu, C. Caramanis, and S. Mannor, ``Robustness and regularization of support vector machines,'' \emph{J. Mach. Learn. Res.}, vol. 10, pp. 1485--1510, Dec. 2009.
		
		\bibitem{muandet2016kernel}
		K. Muandet, B. Sriperumbudur, K. Fukumizu, A. Gretton, and B. Schölkopf, ``Kernel mean shrinkage estimators,'' \emph{J. Mach. Learn. Res.}, vol. 17, no. 1, pp. 1--41, 2016.
		
		\bibitem{flaxman2016bayesian}
		S. Flaxman, D. Sejdinovic, J. P. Cunningham, and S. Filippi, ``Bayesian learning of kernel embeddings,'' 2016. [Online]. Available: https://arxiv.org/pdf/1603.02160.pdf.
		
		\bibitem{dasgupta2008hierarchical}
		S. Dasgupta and D. Hsu, ``Hierarchical sampling for active learning,'' in \emph{Proc. 25th Int. Conf. Mach. Learn.}, 2008, pp. 208--215.
		
		\bibitem{cai2012manifold}
		D. Cai and X. He, ``Manifold adaptive experimental design for text categorization,'' \emph{IEEE Trans. Knowl. Data Eng.}, vol. 24, no. 4, pp. 707--719, Apr. 2012.
		
		\bibitem{lewis1995sequential}
		D. D. Lewis and W. A. Gale, ``A sequential algorithm for training text classifiers,'' \emph{ACM SIGIR Forum}, vol. 29, no. 2, pp. 13--19, 1995.
		
		\bibitem{zhang2015active}
		J. Zhang, X. Wu, and V. S. Sheng, ``Active learning with imbalanced multiple noisy labeling,'' \emph{IEEE Trans. Cybern.}, vol. 45, no. 5, pp. 1095--1107, May 2015.
		
		\bibitem{freund1997selective}
		Y. Freund, H. S. Seung, E. Shamir, and N. Tishby, ``Selective sampling using the query by committee,'' \emph{Mach. Learn.}, vol. 28, pp. 133--168, Aug. 1997.
		
		\bibitem{wang2017committee}
		H. Wang, Y. Jin, and J. Doherty, ``Committee-based active learning for surrogate-assisted particle swarm optimization of expensive problems,'' \emph{IEEE Trans. Cybern.}, vol. 47, no. 9, pp. 2664--2677, Sep. 2017.
		
		\bibitem{tong2001support}
		S. Tong and D. Koller, ``Support vector machine active learning with applications to text classification,'' \emph{J. Mach. Learn. Res.}, vol. 2, pp. 45--66, Mar. 2002.
		
		\bibitem{hoi2008semi}
		S. C. H. Hoi, R. Jin, J. Zhu, and M. R. Lyu, ``Semi-supervised SVM batch mode active learning for image retrieval,'' in \emph{Proc. CVPR}, 2008, pp. 1--7.
		
		\bibitem{huang2014active}
		S.-J. Huang, R. Jin, and Z.-H. Zhou, ``Active learning by querying informative and representative examples,'' \emph{IEEE Trans. Pattern Anal. Mach. Intell.}, vol. 36, no. 10, pp. 1936--1949, Oct. 2014.
		
		\bibitem{du2017exploring}
		B. Du et al., ``Exploring representativeness and informativeness for active learning,'' \emph{IEEE Trans. Cybern.}, vol. 47, no. 1, pp. 14--26, Jan. 2017.
		
		\bibitem{chang2011libsvm}
		C.-C. Chang and C.-J. Lin, ``LIBSVM: A library for support vector machines,'' \emph{ACM Trans. Intell. Syst. Technol.}, vol. 2, no. 3, p. 27, 2011.
		
		\bibitem{fang2017active}
		M. Fang, J. Yin, L. O. Hall, and D. Tao, ``Active multitask learning with trace norm regularization based on excess risk,'' \emph{IEEE Trans. Cybern.}, vol. 47, no. 11, pp. 3906--3915, Nov. 2017.
		
		\bibitem{ramirez2017active}
		M. E. Ramirez-Loaiza, M. Sharma, G. Kumar, and M. Bilgic, ``Active learning: An empirical study of common baselines,'' \emph{Data Min. Knowl. Disc.}, vol. 31, no. 2, pp. 287--313, 2017.
		
		\bibitem{settles2012active}
		B. Settles, \emph{Active Learning} (Synthesis Lectures on Artificial Intellience and Machine Learning). San Rafael, CA, USA: Morgan \& Claypool Publ., 2012.
		
		\bibitem{smola2007hilbert}
		A. Smola, A. Gretton, L. Song, and B. Schölkopf, ``A Hilbert space embedding for distributions,'' in \emph{Proc. Int. Conf. Algorithmic Learn. Theory}. Heidelberg, Germany: Springer, 2007, pp. 13--31.
		
		\bibitem{lopez2015towards}
		D. Lopez-Paz, K. Muandet, B. Schölkopf, and I. O. Tolstikhin, ``Towards a learning theory of cause-effect inference,'' in \emph{Proc. 32nd Int. Conf. Mach. Learn.}, vol. 32, 2015, pp. 1452--1461.
		
		\bibitem{law2017bayesian}
		H. C. L. Law, D. J. Sutherland, D. Sejdinovic, and S. Flaxman, ``Bayesian distribution regression,'' 2017. [Online]. Available: http://arxiv.org/abs/1705.04293.
		
		\bibitem{devroye1996probabilistic}
		L. Devroye, L. Gyorfi, and G. Lugosi, \emph{A Probabilistic Theory of Pattern Recognition}. New York, NY, USA: Springer, 1996.
		
		\bibitem{chen2006consistency}
		D.-R. Chen and T. Sun, ``Consistency of multiclass empirical risk minimization methods based on convex loss,'' \emph{J. Mach. Learn. Res.}, vol. 7, no. 86, pp. 2435--2447, 2006.
		
		\bibitem{sinha2009semi}
		K. Sinha and M. Belkin, ``Semi-supervised learning using sparse eigenfunction bases,'' in \emph{Advances in Neural Information Processing Systems 22}. La Jolla, CA, USA: Curran Assoc., Inc., 2009, pp. 1687--1695.
		
		\bibitem{singh2009unlabeled}
		A. Singh and R. D. Nowak, ``Unlabeled data: Now it helps, now it doesn't,'' in \emph{Advances in Neural Information Processing Systems 21}. La Jolla, CA, USA: Curran Assoc., Inc., 2009, pp. 1513--1520.
		
		\bibitem{audibert2007fast}
		J. Y. Audibert and A. B. Tsybakov, ``Fast learning rates for plug-in classifiers,'' \emph{Ann. Stat.}, vol. 35, no. 2, pp. 608--633, 2007.
		
		\bibitem{rasmussen2004gaussian}
		C. E. Rasmussen, ``Gaussian processes in machine learning,'' in \emph{Advanced Lectures on Machine Learning}, vol. 3176. Heidelberg, Germany: Springer, 2004, pp. 63--71.
		
		\bibitem{goemans1995improved}
		M. X. Goemans and D. P. Williamson, ``Improved approximation algorithms for maximum cut and satisfiability problems using semidefinite programming,'' \emph{J. Assoc. Comput. Mach.}, vol. 42, no. 6, pp. 1115--1145, 1995.
		
		\bibitem{wong2001asymptotic}
		R. Wong, \emph{Asymptotic Approximations of Integrals}. Philadelphia, PA, USA: SIAM, 2001.
		
		\bibitem{lapinski2019multivariate}
		T. M. Łapiński, ``Multivariate Laplace's approximation with estimated error and application to limit theorems,'' \emph{J. Approx. Theory}, vol. 248, Dec. 2019, Art. no. 105305.
		
		\bibitem{inglot2014simple}
		T. Inglot and P. Majerski, ``Simple upper and lower bounds for the multivariate Laplace approximation,'' \emph{J. Approx. Theory}, vol. 186, pp. 1--11, Oct. 2014.
		
		\bibitem{zhu2003semi}
		X. Zhu, Z. Ghahramani, and J. Lafferty, ``Combining active learning and semi-supervised learning using Gaussian fields and harmonic functions,'' in \emph{Proc. Int. Conf. Mach. Learn.}, vol. 20, 2003, pp. 912--919.
		
		\bibitem{mikolov2013efficient}
		T. Mikolov, G. Corrado, K. Chen, and J. Dean, ``Efficient estimation of word representations in vector space,'' in \emph{Proc. Int. Conf. Learn. Represent. (ICLR)}, 2013, pp. 1--12.
		
		\bibitem{nene1996columbia}
		S. Nene, S. Nayar, and H. Murase, ``Columbia object image library (COIL-20),'' Dept. Comput. Sci., Columbia Univ., New York, NY, USA, Rep. CUCS-005-96, 1996.
		
		\bibitem{teufel2002summarizing}
		S. Teufel and M. Moens, ``Summarizing scientific articles: Experiments with relevance and rhetorical status,'' \emph{Comput. Linguist.}, vol. 28, no. 4, pp. 409--445, 2002.
		
	\end{thebibliography}
	
\end{document}